\documentclass[12pt, a4paper]{article}
% --- 1. Cấu hình Tiếng Việt và Font chữ chuẩn ---
\usepackage[utf8]{inputenc}
\usepackage[T5]{fontenc}
\usepackage[vietnamese]{babel}
\usepackage{mathptmx} % Font Times New Roman đồng bộ cả chữ và công thức toán, không gây xung đột gói

\makeatletter
\renewcommand{\normalsize}{\@setfontsize\normalsize{13pt}{16pt}}
\makeatother
\normalsize

% --- 2. Gói Toán học & Ký hiệu bổ trợ ---
\usepackage{amsmath, amssymb, amsfonts, mathrsfs, amsthm}
\usepackage{graphicx}
\usepackage{fontawesome}
\usepackage{booktabs, tabularx, array, multirow}
\usepackage{enumitem}

% --- 3. Đồ họa TikZ, Bảng biến thiên & Hình không gian ---
\usepackage{tikz}
\usetikzlibrary{calc, intersections, angles, quotes, patterns, arrows.meta, 3d}
\usepackage{pgfplots}
\pgfplotsset{compat=1.18}
\usepackage{tkz-euclide}
\usepackage{tkz-tab}

\renewcommand{\vec}[1]{\overrightarrow{#1}}

% --- 4. Hộp sư phạm (Tcolorbox) ---
\usepackage[many]{tcolorbox}
\tcbuselibrary{skins, breakable}

% --- 5. Căn lề chuẩn văn bản giáo án ---
\usepackage{geometry}
\geometry{
    a4paper,
    left=25mm,
    right=15mm,
    top=20mm,
    bottom=20mm
}

% --- 6. Header / Footer chuẩn hóa ---
\usepackage{fancyhdr}
\usepackage{lastpage}
\pagestyle{fancy}
\fancyhf{}

\fancyhead[L]{\small \textit{Kế hoạch bài dạy Toán 11}}
\fancyhead[R]{\textbf{Trang \thepage/\pageref{LastPage}}}
\fancyfoot[L]{\small \textit{Tổ Toán - Tin}}
\fancyfoot[R]{\small \textit{Trường THCS\&THPT Hồng Vân}}
\renewcommand{\headrulewidth}{0.4pt}
\renewcommand{\footrulewidth}{0.4pt}

% --- 7. Giãn dòng & Giãn đoạn theo chuẩn Nghị định 30/2020/NĐ-CP ---
\usepackage{setspace}
\setstretch{1.15}
\setlength{\parindent}{1.0cm}
\setlength{\parskip}{5pt}

% Định nghĩa hộp sư phạm chuyên dụng
\newtcolorbox{pedabox}[2][]{
    enhanced,
    breakable,
    colback=blue!3!white,
    colframe=blue!65!black,
    fonttitle=\bfseries\small,
    title={#2},
    arc=2mm,
    boxrule=0.6pt,
    left=3mm, right=3mm, top=2mm, bottom=2mm,
    #1
}

\newtcolorbox{aibox}[2][]{
    enhanced,
    breakable,
    colback=teal!4!white,
    colframe=teal!70!black,
    fonttitle=\bfseries\small,
    title={#2},
    arc=2mm,
    boxrule=0.6pt,
    left=3mm, right=3mm, top=2mm, bottom=2mm,
    #1
}

\newtcolorbox{scaffoldbox}[2][]{
    enhanced,
    breakable,
    colback=orange!4!white,
    colframe=orange!80!black,
    fonttitle=\bfseries\small,
    title={#2},
    arc=2mm,
    boxrule=0.6pt,
    left=3mm, right=3mm, top=2mm, bottom=2mm,
    #1
}

\begin{document}

\begin{center}
    {\fontsize{14pt}{17pt}\selectfont \textbf{KẾ HOẠCH BÀI DẠY MÔN TOÁN LỚP 11}}\\[6pt]
    {\fontsize{16pt}{19pt}\selectfont \textbf{BÀI 3. HÀM SỐ LƯỢNG GIÁC}}\\[4pt]
    {\fontsize{12pt}{15pt}\selectfont \textbf{Môn học: Toán 11} \ \textbar \ \textbf{Thời lượng: 2 tiết (Tiết 6 và Tiết 7 theo PPCT)}}
\end{center}

\vspace{5pt}

\section*{I. MỤC TIÊU}

\subsection*{1. Kiến thức, Năng lực}

\begin{itemize}[leftmargin=1.2cm]
    \item \textbf{Yêu cầu cần đạt (Nguyên văn CT GDPT 2018):}
    \begin{itemize}[leftmargin=0.6cm]
        \item Nhận biết được các khái niệm về hàm số chẵn, hàm số lẻ, hàm số tuần hoàn.
        \item Nhận biết được các đặc trưng hình học của đồ thị hàm số chẵn, hàm số lẻ, hàm số tuần hoàn.
        \item Nhận biết được định nghĩa các hàm lượng giác $y = \sin x, y = \cos x, y = \tan x, y = \cot x$ thông qua đường tròn lượng giác.
        \item Mô tả được bảng giá trị của bốn hàm số lượng giác đó trên một chu kì.
        \item Vẽ được đồ thị của các hàm số $y = \sin x, y = \cos x, y = \tan x, y = \cot x$.
        \item Giải thích được: tập xác định; tập giá trị; tính chất chẵn, lẻ; tính tuần hoàn; chu kì; khoảng đồng biến, nghịch biến của các hàm số lượng giác dựa vào đồ thị.
        \item Giải quyết được một số vấn đề thực tiễn gắn với hàm số lượng giác (ví dụ: dao động điều hoà trong Vật lí,...).
    \end{itemize}

    \item \textbf{Năng lực Toán học đặc thù:}
    \begin{itemize}[leftmargin=0.6cm]
        \item \textit{Năng lực tư duy và lập luận toán học:} So sánh sự khác nhau về tính tuần hoàn, tập xác định và tập giá trị giữa hàm $\sin, \cos$ và hàm $\tan, \cot$; phân tích trục đối xứng ($Oy$) của đồ thị hàm số chẵn và tâm đối xứng (gốc $O$) của đồ thị hàm số lẻ; suy luận khoảng đồng biến, nghịch biến của các hàm số lượng giác từ đồ thị.
        \item \textit{Năng lực mô hình hóa toán học:} Thiết lập hàm số lượng giác dạng $y = A \cos(\omega t + \varphi)$ hoặc $y = A \sin(\omega t + \varphi)$ để mô tả các hiện tượng biến đổi tuần hoàn trong tự nhiên và đời sống: chuyển động của con lắc lò xo, đồ thị dao động điều hòa, quy luật mực nước thủy triều lên xuống theo thời gian.
        \item \textit{Năng lực giải quyết vấn đề toán học:} Tìm tập xác định của các hàm số lượng giác chứa phân thức, căn thức; xác định giá trị lớn nhất, giá trị nhỏ nhất của hàm số lượng giác đơn giản; tìm chu kỳ tuần hoàn.
        \item \textit{Năng lực giao tiếp toán học:} Trình bày rõ ràng các bước tìm tập xác định, xét tính chẵn lẻ và phân tích sự biến thiên của hàm số; sử dụng chuẩn xác các ký hiệu toán học $D = \mathbb{R} \setminus \{\dots\}$.
        \item \textit{Năng lực sử dụng công cụ, phương tiện học toán:} Khai thác máy tính cầm tay (MTCT) Casio fx-580VN X bằng bảng giá trị (\texttt{TABLE}) để khảo sát quy luật tuần hoàn; sử dụng phần mềm hình học động GeoGebra để vẽ và quan sát dạng sóng hình sin của đồ thị hàm số lượng giác.
    \end{itemize}

    \item \textbf{Năng lực số (Theo Thông tư 02/2025/TT-BGDĐT):}
    \begin{itemize}[leftmargin=0.6cm]
        \item \textit{Mã 3.1NC1b:} Học sinh quan sát, phân tích đồ thị hàm số $y = \sin x, y = \cos x$ trên màn hình trình chiếu và trực tiếp thực hành vẽ đồ thị động trên phần mềm GeoGebra tại phòng máy vi tính để nhận diện trực quan chu kì tuần hoàn $2\pi$ và $\pi$.
    \end{itemize}

    \item \textbf{Năng lực AI (Theo Quyết định 2422/QĐ-BGDĐT):}
    \begin{itemize}[leftmargin=0.6cm]
        \item \textit{Mã 11.C3.2:} Học sinh mô tả và trình bày được cách công nghệ Trí tuệ nhân tạo (AI) nhận dạng, phân tích và dự báo các chuỗi dữ liệu có tính chu kì trong y tế và công nghệ (sóng âm thanh, tín hiệu điện tâm đồ ECG, nhịp sinh học).
        \item \textit{Mã 11.C3.1:} Học sinh thực hành kỹ thuật đặt câu lệnh (prompt) chuẩn mực yêu cầu trợ lý AI phân tích các khoảng đồng biến, nghịch biến của hàm số lượng giác trên một chu kì, từ đó so sánh và kiểm chứng với bài giải giải tích.
    \end{itemize}

    \item \textbf{Năng lực chung:}
    \begin{itemize}[leftmargin=0.6cm]
        \item \textit{Tự chủ và tự học:} Tự giác nghiên cứu tài liệu SGK, hoàn thành nhiệm vụ khảo sát hàm số trên phiếu học tập.
        \item \textit{Giao tiếp và hợp tác:} Tương tác tích cực trong thảo luận nhóm, phân công vẽ đồ thị và đối chiếu kết quả.
    \end{itemize}
\end{itemize}

\subsection*{2. Phẩm chất}

\begin{itemize}[leftmargin=1.2cm]
    \item \textbf{Chăm chỉ:} Tỉ mỉ lập bảng giá trị, kiên trì vẽ đồ thị uốn lượn chính xác, tích cực thực hành phần mềm số.
    \item \textbf{Trung thực:} Báo cáo chính xác kết quả khảo sát hàm số, trung thực trong kiểm tra chéo và đối chiếu với AI.
    \item \textbf{Trách nhiệm:} Tích cực hỗ trợ bạn yếu trong việc xác định điều kiện mẫu số khác 0 của hàm $\tan, \cot$; có ý thức ứng dụng toán học vào phân tích hiện tượng tự nhiên.
\end{itemize}

\section*{II. THIẾT BỊ DẠY HỌC VÀ HỌC LIỆU}

\begin{itemize}[leftmargin=1.2cm]
    \item \textbf{Giáo viên:}
    \begin{itemize}[leftmargin=0.6cm]
        \item Kế hoạch bài dạy chi tiết, bài trình chiếu đa phương tiện (PowerPoint/Canva) tích hợp đồ thị động.
        \item File GeoGebra mô phỏng chuyển động của điểm $M$ trên đường tròn lượng giác và vết để lại tạo thành đồ thị hình sin.
        \item Video ngắn minh họa ứng dụng sóng âm thanh, tín hiệu nhịp tim điện tâm đồ ECG và chuyển động thủy triều.
        \item Hệ thống Phiếu học tập phân tầng (Worksheet 1, Worksheet 2) và Rubric đánh giá năng lực 4 mức độ.
    \end{itemize}
    \item \textbf{Học sinh:}
    \begin{itemize}[leftmargin=0.6cm]
        \item SGK Toán 11, vở ghi bài, giấy kẻ ô li để vẽ đồ thị.
        \item Thước kẻ parabol/thước dẻo, bút màu, MTCT Casio fx-580VN X.
        \item Điện thoại thông minh/máy tính phòng thực hành có kết nối Internet để chạy GeoGebra.
    \end{itemize}
\end{itemize}

\section*{III. TIẾN TRÌNH DẠY HỌC}

% =========================================================================
% TIẾT 1 CỦA BÀI 3 (TIẾT 6 PPCT)
% =========================================================================
\subsection*{TIẾT 1. HÀM SỐ TUẦN HOÀN, CHẴN LẺ VÀ HÀM SỐ SIN, COSIN}
\textit{(Tiết 6 theo PPCT | Tuần 2)}

\subsubsection*{1. Hoạt động 1: Khởi động (Mở đầu)}
\textbf{a) Mục tiêu:} Tạo tình huống thực tế về các quá trình lặp đi lặp lại có chu kỳ trong tự nhiên (nhịp tim, cánh quạt quay, tiếng còi xe), khơi gợi nhu cầu định nghĩa Hàm số tuần hoàn.

\textbf{b) Nội dung:} Giáo viên chiếu video đồ thị máy đo điện tâm đồ (ECG) hiển thị nhịp tim đập đều đặn và đồ thị mực nước thủy triều theo các giờ trong ngày. Học sinh quan sát và trả lời câu hỏi:
\begin{itemize}[leftmargin=0.6cm]
    \item Hình dạng sóng trên màn hình có đặc điểm gì nổi bật sau những khoảng thời gian bằng nhau?
    \item Trong các hàm số đã học ở lớp 10 (hàm bậc nhất $y = ax + b$, hàm bậc hai $y = ax^2 + bx + c$), có hàm số nào có tính chất lặp lại sau một khoảng cách cố định như vậy không?
\end{itemize}

\textbf{c) Sản phẩm:} Học sinh trả lời:
\begin{itemize}[leftmargin=0.6cm]
    \item Đồ thị nhịp tim và thủy triều có dạng lượn sóng lặp đi lặp lại giống hệt nhau sau một khoảng thời gian xác định $T$.
    \item Các hàm đa thức lớp 10 không có tính chất này (chúng tiến ra vô cực khi $x$ tăng). Cần một lớp hàm số mới gọi là \textbf{Hàm số tuần hoàn}.
\end{itemize}

\textbf{d) Tổ chức thực hiện:}
\begin{itemize}[leftmargin=0.6cm]
    \item \textit{Chuyển giao nhiệm vụ:} Giáo viên trình chiếu video, đặt câu hỏi cho học sinh suy nghĩ trong 2 phút.
    \item \textit{Thực hiện nhiệm vụ (Giàn giáo sư phạm):}
    \begin{scaffoldbox}{Giàn giáo sư phạm 3.1: Nhận diện hiện tượng tuần hoàn}
        \textbf{Câu hỏi dẫn dắt:} Khi kim giây đồng hồ quay, cứ sau bao nhiêu giây thì nó trở về đúng vị trí ban đầu? (Trả lời: 60 giây). Vậy các giá trị vị trí lặp lại sau một chu kỳ cố định.
    \end{scaffoldbox}
    \item \textit{Báo cáo, thảo luận:} 2 học sinh phát biểu, cả lớp nhận xét.
    \item \textit{Kết luận, nhận định:} Giáo viên dẫn dắt vào bài mới: Thiết lập định nghĩa Hàm số tuần hoàn và nghiên cứu hai hàm số lượng giác đầu tiên: $y = \sin x$ và $y = \cos x$.
\end{itemize}

\subsubsection*{2. Hoạt động 2: Hình thành kiến thức}

\paragraph{Nội dung 1: Khái niệm hàm số tuần hoàn, hàm số chẵn, hàm số lẻ}
\textbf{a) Mục tiêu:} Nhận biết định nghĩa hàm số tuần hoàn và chu kỳ tuần hoàn; củng cố định nghĩa và đặc trưng đối xứng hình học của hàm số chẵn, hàm số lẻ.

\textbf{b) Nội dung:}
\begin{itemize}[leftmargin=0.6cm]
    \item \textbf{Định nghĩa hàm số tuần hoàn:} Hàm số $y = f(x)$ có tập xác định $D$ được gọi là tuần hoàn nếu tồn tại một số $T \neq 0$ sao cho với mọi $x \in D$:
    \begin{enumerate}[leftmargin=0.6cm]
        \item $x - T \in D$ và $x + T \in D$.
        \item $f(x + T) = f(x)$.
    \end{enumerate}
    Số dương $T$ nhỏ nhất thỏa mãn các điều kiện trên (nếu có) được gọi là \textbf{chu kì} của hàm số tuần hoàn đó.
    \item \textbf{Hàm số chẵn và hàm số lẻ (Đặc trưng hình học đồ thị):}
    \begin{itemize}[leftmargin=0.5cm]
        \item $f(x)$ là hàm số chẵn $\iff \forall x \in D \implies -x \in D$ và $f(-x) = f(x)$.\\
        $\implies$ Đồ thị của hàm số chẵn nhận \textbf{trục tung $Oy$ làm trục đối xứng}.
        \item $f(x)$ là hàm số lẻ $\iff \forall x \in D \implies -x \in D$ và $f(-x) = -f(x)$.\\
        $\implies$ Đồ thị của hàm số lẻ nhận \textbf{gốc toạ độ $O$ làm tâm đối xứng}.
    \end{itemize}
\end{itemize}

\textbf{c) Sản phẩm:} Học sinh phân biệt được hàm chẵn và hàm lẻ qua công thức lượng giác:
\begin{itemize}[leftmargin=0.6cm]
    \item $\cos(-x) = \cos x \implies y = \cos x$ là hàm số chẵn.
    \item $\sin(-x) = -\sin x \implies y = \sin x$ là hàm số lẻ.
\end{itemize}

\textbf{d) Tổ chức thực hiện:}
\begin{itemize}[leftmargin=0.6cm]
    \item \textit{Chuyển giao nhiệm vụ:} Giáo viên trình bày định nghĩa, yêu cầu học sinh làm việc cá nhân nhắc lại tính chất đối xứng của đồ thị hàm chẵn/lẻ.
    \item \textit{Thực hiện nhiệm vụ (Giàn giáo cho HS TB-Yếu):}
    \begin{scaffoldbox}{Giàn giáo sư phạm 3.2: Mẹo nhớ trục đối xứng và tâm đối xứng}
        \textbf{Khắc sâu trực quan:}\\
        - Chữ \textbf{CHẴN} có trục đối xứng là $Oy$ (giống hình ảnh chiếc bình hoa đối xứng qua gương đứng).\\
        - Chữ \textbf{LẺ} có tâm đối xứng là gốc tọa độ $O$ (giống hình cánh quạt 2 cánh quay $180^\circ$ quanh tâm).
    \end{scaffoldbox}
    \item \textit{Báo cáo, thảo luận:} 1 học sinh trả lời miệng, giáo viên chuẩn hóa.
    \item \textit{Kết luận, nhận định:} Giáo viên nhấn mạnh: Tính chất đối xứng giúp ta chỉ cần khảo sát đồ thị trên một nửa miền xác định rồi lấy đối xứng.
\end{itemize}

\paragraph{Nội dung 2: Định nghĩa và tính chất của hàm số $y = \sin x$ và $y = \cos x$}
\textbf{a) Mục tiêu:} Nhận biết định nghĩa hàm số $y = \sin x$ và $y = \cos x$ thông qua đường tròn lượng giác; xác định tập xác định, tập giá trị, tính chẵn lẻ, tính tuần hoàn và chu kỳ $2\pi$.

\textbf{b) Nội dung:}
\begin{itemize}[leftmargin=0.6cm]
    \item \textbf{Hàm số sin ($y = \sin x$):}
    \begin{itemize}[leftmargin=0.5cm]
        \item Định nghĩa: Quy tắc đặt tương ứng mỗi số thực $x$ với số thực $\sin x$ (tung độ điểm biểu diễn góc $x$ trên đường tròn lượng giác).
        \item Tập xác định: $D = \mathbb{R}$.
        \item Tập giá trị: $T = [-1; 1]$, nghĩa là $-1 \le \sin x \le 1$ với mọi $x \in \mathbb{R}$.
        \item Tính chẵn lẻ: Là hàm số \textbf{lẻ} vì $\sin(-x) = -\sin x$.
        \item Tính tuần hoàn: Tuần hoàn với chu kì \textbf{$T = 2\pi$} vì $\sin(x + 2\pi) = \sin x$.
    \end{itemize}
    \item \textbf{Hàm số cosin ($y = \cos x$):}
    \begin{itemize}[leftmargin=0.5cm]
        \item Định nghĩa: Quy tắc đặt tương ứng mỗi số thực $x$ với số thực $\cos x$ (hoành độ điểm biểu diễn góc $x$).
        \item Tập xác định: $D = \mathbb{R}$.
        \item Tập giá trị: $T = [-1; 1]$, nghĩa là $-1 \le \cos x \le 1$ với mọi $x \in \mathbb{R}$.
        \item Tính chẵn lẻ: Là hàm số \textbf{chẵn} vì $\cos(-x) = \cos x$.
        \item Tính tuần hoàn: Tuần hoàn với chu kì \textbf{$T = 2\pi$} vì $\cos(x + 2\pi) = \cos x$.
    \end{itemize}
\end{itemize}

\textbf{c) Sản phẩm:} Bảng so sánh đặc trưng của 2 hàm số:
\begin{center}
\begin{tabular}{lcc}
\toprule
\textbf{Đặc trưng} & \textbf{Hàm số $y = \sin x$} & \textbf{Hàm số $y = \cos x$} \\
\midrule
Tập xác định & $D = \mathbb{R}$ & $D = \mathbb{R}$ \\
Tập giá trị & $[-1; 1]$ & $[-1; 1]$ \\
Tính chất chẵn / lẻ & Hàm số lẻ (đối xứng qua gốc $O$) & Hàm số chẵn (đối xứng qua trục $Oy$) \\
Chu kì tuần hoàn & $T = 2\pi$ & $T = 2\pi$ \\
\bottomrule
\end{tabular}
\end{center}

\textbf{d) Tổ chức thực hiện (Tích hợp Năng lực số 3.1NC1b \& Năng lực AI 11.C3.2):}
\begin{itemize}[leftmargin=0.6cm]
    \item \textit{Chuyển giao nhiệm vụ:} Giáo viên trình chiếu mô phỏng GeoGebra điểm chạy trên đường tròn đơn vị và chiếu vết tung độ ra trục số.
    \item \textit{Thực hiện nhiệm vụ (Năng lực số):} Học sinh quan sát màn hình, thấy khi góc $x$ tăng từ $0$ lên $2\pi$ thì tung độ $\sin x$ đi từ $0 \to 1 \to 0 \to -1 \to 0$ và chu trình lặp lại y hệt khi vượt qua $2\pi$.
    \item \textit{Tích hợp Năng lực AI (Mã 11.C3.2):}
    \begin{aibox}{Thực hành Năng lực AI: Nhận diện dữ liệu chu kỳ trong AI}
        \textbf{Ứng dụng thực tế:}\\
        Trong y tế, mô hình học sâu (Deep Learning AI) sử dụng các hàm lượng giác chu kỳ để phân tích dữ liệu điện tâm đồ (ECG) và tín hiệu sóng não (EEG). Bằng cách so sánh chu kỳ thực tế với chu kỳ chuẩn $T$, AI có thể phát hiện sớm hiện tượng loạn nhịp tim với độ chính xác trên 98\%.
    \end{aibox}
    \item \textit{Báo cáo, thảo luận:} 1 học sinh đại diện nhắc lại lý do tại sao chu kỳ của $\sin x$ và $\cos x$ là $2\pi$ chứ không phải $\pi$.
    \item \textit{Kết luận, nhận định:} Giáo viên chốt kiến thức: Vì sau đúng $2\pi$ (1 vòng tròn), điểm biểu diễn mới quay trở lại vị trí cũ.
\end{itemize}

\subsubsection*{3. Hoạt động 3: Luyện tập}
\textbf{a) Mục tiêu:} Rèn luyện kỹ năng tìm tập xác định, xét tính chẵn lẻ của hàm số chứa sin và cosin; tìm giá trị lớn nhất, nhỏ nhất.

\textbf{b) Nội dung:} Học sinh hoàn thành Phiếu học tập số 1 (Worksheet 1):
\begin{itemize}[leftmargin=0.6cm]
    \item \textbf{Bài 1:} Tìm tập xác định của các hàm số sau:
    $$a) \ y = \dfrac{1}{\sin x}; \quad b) \ y = \sqrt{1 - \cos x}$$
    \item \textbf{Bài 2:} Xét tính chẵn, lẻ của các hàm số sau:
    $$a) \ f(x) = x \sin x; \quad b) \ g(x) = \sin x + \cos x$$
    \item \textbf{Bài 3:} Tìm giá trị lớn nhất ($M$) và giá trị nhỏ nhất ($m$) của hàm số:
    $$y = 3\cos x - 2$$
\end{itemize}

\textbf{c) Sản phẩm:} Lời giải chi tiết:
\begin{itemize}[leftmargin=0.6cm]
    \item \textbf{Bài 1:}
    \begin{itemize}[leftmargin=0.5cm]
        \item a) Hàm số xác định $\iff \sin x \neq 0 \iff x \neq k\pi$ ($k \in \mathbb{Z}$). Vậy $D = \mathbb{R} \setminus \{k\pi, k \in \mathbb{Z}\}$.
        \item b) Vì $-1 \le \cos x \le 1 \implies 1 - \cos x \ge 0$ với mọi $x \in \mathbb{R}$. Do đó hàm số xác định với mọi $x \in \mathbb{R} \implies D = \mathbb{R}$.
    \end{itemize}
    \item \textbf{Bài 2:}
    \begin{itemize}[leftmargin=0.5cm]
        \item a) Tập xác định $D = \mathbb{R}$. Với mọi $x \in \mathbb{R} \implies -x \in \mathbb{R}$.\\
        $f(-x) = (-x) \sin(-x) = (-x)(-\sin x) = x \sin x = f(x) \implies f(x)$ là \textbf{hàm số chẵn}.
        \item b) $D = \mathbb{R}$. Ta có $g(-x) = \sin(-x) + \cos(-x) = -\sin x + \cos x$.\\
        Xét $g\left(\dfrac{\pi}{4}\right) = \dfrac{\sqrt{2}}{2} + \dfrac{\sqrt{2}}{2} = \sqrt{2}$, nhưng $g\left(-\dfrac{\pi}{4}\right) = -\dfrac{\sqrt{2}}{2} + \dfrac{\sqrt{2}}{2} = 0 \neq \pm \sqrt{2}$.\\
        Vậy $g(x)$ \textbf{không phải là hàm số chẵn và cũng không phải là hàm số lẻ}.
    \end{itemize}
    \item \textbf{Bài 3:}
    Vì $-1 \le \cos x \le 1 \implies -3 \le 3\cos x \le 3 \implies -3 - 2 \le 3\cos x - 2 \le 3 - 2 \iff -5 \le y \le 1$.\\
    Vậy $\max y = 1$ (khi $\cos x = 1 \iff x = k2\pi$); $\min y = -5$ (khi $\cos x = -1 \iff x = \pi + k2\pi$).
\end{itemize}

\textbf{d) Tổ chức thực hiện:}
\begin{itemize}[leftmargin=0.6cm]
    \item \textit{Chuyển giao nhiệm vụ:} Học sinh làm bài cá nhân 7 phút, đổi chéo chấm điểm theo cặp.
    \item \textit{Thực hiện nhiệm vụ (Hỗ trợ HS TB-Yếu):} Giáo viên nhắc học sinh yếu: Hàm số không chẵn không lẻ là trường hợp rất phổ biến, không được kết luận vội vàng nếu không thỏa mãn cả 2 điều kiện.
    \item \textit{Báo cáo, thảo luận:} 3 học sinh lên bảng trình bày, cả lớp đối chiếu.
    \item \textit{Kết luận, nhận định:} Giáo viên chuẩn hóa cách trình bày tìm tập xác định và tập giá trị.
\end{itemize}

\subsubsection*{4. Hoạt động 4: Vận dụng}
\textbf{a) Mục tiêu:} Vận dụng hàm số cosin để mô tả quy luật dao động điều hòa của con lắc lò xo.

\textbf{b) Nội dung:}
Một con lắc lò xo dao động điều hòa với phương trình li độ: $x(t) = 4\cos\left(2\pi t - \dfrac{\pi}{3}\right)$ ($x$ tính bằng cm, $t$ tính bằng giây).
\begin{enumerate}[leftmargin=0.5cm]
    \item Tìm biên độ dao động $A$ và chu kỳ dao động $T$ của con lắc.
    \item Xác định vị trí của con lắc tại thời điểm ban đầu $t = 0\text{ s}$ và thời điểm $t = 0{,}5\text{ s}$.
\end{enumerate}

\textbf{c) Sản phẩm:} Lời giải:
\begin{enumerate}[leftmargin=0.5cm]
    \item Biên độ dao động $A = 4\text{ cm}$.\\
    Chu kỳ dao động $T = \dfrac{2\pi}{\omega} = \dfrac{2\pi}{2\pi} = 1\text{ s}$.
    \item Tại $t = 0$: $x(0) = 4\cos\left(-\dfrac{\pi}{3}\right) = 4 \cdot \dfrac{1}{2} = 2\text{ cm}$.\\
    Tại $t = 0{,}5$: $x(0{,}5) = 4\cos\left(2\pi \cdot 0{,}5 - \dfrac{\pi}{3}\right) = 4\cos\left(\pi - \dfrac{\pi}{3}\right) = 4\cos\dfrac{2\pi}{3} = 4 \cdot \left(-\dfrac{1}{2}\right) = -2\text{ cm}$.
\end{enumerate}

\textbf{d) Tổ chức thực hiện:}
\begin{itemize}[leftmargin=0.6cm]
    \item \textit{Chuyển giao nhiệm vụ:} Giáo viên giao bài toán liên môn Vật lý 11 cho cả lớp.
    \item \textit{Thực hiện nhiệm vụ:} Học sinh tính toán nhanh bằng MTCT trong 3 phút.
    \item \textit{Báo cáo, thảo luận:} 1 học sinh trả lời tại chỗ, giải thích ý nghĩa dấu âm của li độ.
    \item \textit{Kết luận, nhận định:} Giáo viên tổng kết Tiết 1, dặn dò học sinh xem trước đồ thị của hàm số tang và cotang cho Tiết 2.
\end{itemize}

% =========================================================================
% TIẾT 2 CỦA BÀI 3 (TIẾT 7 PPCT)
% =========================================================================
\vspace{10pt}
\subsection*{TIẾT 2. HÀM SỐ TANG, COTANG VÀ ĐỒ THỊ CỦA CÁC HÀM SỐ LƯỢNG GIÁC}
\textit{(Tiết 7 theo PPCT | Tuần 3)}

\subsubsection*{1. Hoạt động 1: Khởi động (Mở đầu)}
\textbf{a) Mục tiêu:} Khám phá sự khác biệt về tập xác định giữa hàm $\tan, \cot$ (có mẫu số) và hàm $\sin, \cos$; dẫn dắt vào tính chất chu kỳ $\pi$.

\textbf{b) Nội dung:} Giáo viên yêu cầu học sinh dùng máy tính cầm tay bấm các giá trị sau:
$$\tan 89^\circ; \quad \tan 89{,}9^\circ; \quad \tan 89{,}999^\circ; \quad \tan 90^\circ$$
Quan sát kết quả trên màn hình và nhận xét: Điều gì xảy ra khi góc tiến dần đến $90^\circ$ ($\dfrac{\pi}{2}$)? Tại sao máy báo \texttt{Math ERROR} tại $90^\circ$?

\textbf{c) Sản phẩm:} Học sinh nhận thấy:
\begin{itemize}[leftmargin=0.6cm]
    \item $\tan 89^\circ \approx 57{,}29$; $\tan 89{,}9^\circ \approx 572{,}9$; $\tan 89{,}999^\circ \approx 57295{,}7 \implies$ giá trị tăng vọt ra dương vô cực.
    \item Tại $90^\circ$, máy báo lỗi vì $\cos 90^\circ = 0$, mẫu số bằng 0. Vậy hàm $\tan$ và $\cot$ không xác định trên toàn bộ $\mathbb{R}$.
\end{itemize}

\textbf{d) Tổ chức thực hiện:}
\begin{itemize}[leftmargin=0.6cm]
    \item \textit{Chuyển giao nhiệm vụ:} Học sinh bấm máy tại chỗ trên 35 MTCT trong 1 phút.
    \item \textit{Thực hiện nhiệm vụ:} Cả lớp nhận thấy hiện tượng giá trị tăng vọt và báo lỗi.
    \item \textit{Báo cáo, thảo luận:} 1 học sinh trả lời nguyên nhân do điều kiện $\cos x \neq 0$.
    \item \textit{Kết luận, nhận định:} Giáo viên dẫn dắt vào bài mới: Nghiên cứu tập xác định, chu kỳ và đồ thị của hàm $\tan, \cot$, sau đó tổng hợp trọn vẹn 4 đồ thị lượng giác.
\end{itemize}

\subsubsection*{2. Hoạt động 2: Hình thành kiến thức}

\paragraph{Nội dung 1: Hàm số $y = \tan x$ và $y = \cot x$}
\textbf{a) Mục tiêu:} Nắm vững định nghĩa, tập xác định, tập giá trị, tính chẵn lẻ và tính tuần hoàn với chu kì $\pi$ của hàm số $y = \tan x$ và $y = \cot x$.

\textbf{b) Nội dung:}
\begin{itemize}[leftmargin=0.6cm]
    \item \textbf{Hàm số tang ($y = \tan x$):}
    \begin{itemize}[leftmargin=0.5cm]
        \item Định nghĩa: $y = \tan x = \dfrac{\sin x}{\cos x}$.
        \item Tập xác định: $D = \mathbb{R} \setminus \left\{\dfrac{\pi}{2} + k\pi, k \in \mathbb{Z}\right\}$.
        \item Tập giá trị: $T = \mathbb{R}$.
        \item Tính chẵn lẻ: Là hàm số \textbf{lẻ} vì $\tan(-x) = -\tan x$.
        \item Tính tuần hoàn: Tuần hoàn với chu kì \textbf{$T = \pi$} vì $\tan(x + \pi) = \tan x$.
    \end{itemize}
    \item \textbf{Hàm số cotang ($y = \cot x$):}
    \begin{itemize}[leftmargin=0.5cm]
        \item Định nghĩa: $y = \cot x = \dfrac{\cos x}{\sin x}$.
        \item Tập xác định: $D = \mathbb{R} \setminus \{k\pi, k \in \mathbb{Z}\}$.
        \item Tập giá trị: $T = \mathbb{R}$.
        \item Tính chẵn lẻ: Là hàm số \textbf{lẻ} vì $\cot(-x) = -\cot x$.
        \item Tính tuần hoàn: Tuần hoàn với chu kì \textbf{$T = \pi$} vì $\cot(x + \pi) = \cot x$.
    \end{itemize}
\end{itemize}

\textbf{c) Sản phẩm:} Học sinh ghi chép đầy đủ các đặc trưng và giải thích được tại sao chu kỳ của $\tan$ và $\cot$ là $\pi$ (do công thức hơn kém $\pi$: $\tan(x + \pi) = \tan x$).

\textbf{d) Tổ chức thực hiện:}
\begin{itemize}[leftmargin=0.6cm]
    \item \textit{Chuyển giao nhiệm vụ:} Giáo viên yêu cầu học sinh thảo luận cặp đôi so sánh chu kỳ của $\sin, \cos$ ($2\pi$) với chu kỳ của $\tan, \cot$ ($\pi$).
    \item \textit{Thực hiện nhiệm vụ (Giàn giáo cho HS TB-Yếu):}
    \begin{scaffoldbox}{Giàn giáo sư phạm 3.3: Khắc sâu tập xác định của tan và cot}
        \textbf{Mẹo ghi nhớ điều kiện mẫu số:}
        \begin{itemize}[leftmargin=0.5cm]
            \item $\tan x = \dfrac{\sin x}{\cos x} \implies \cos x \neq 0 \implies x \neq \dfrac{\pi}{2} + k\pi$ (bỏ các điểm trên trục tung $Oy$).
            \item $\cot x = \dfrac{\cos x}{\sin x} \implies \sin x \neq 0 \implies x \neq k\pi$ (bỏ các điểm trên trục hoành $Ox$).
        \end{itemize}
    \end{scaffoldbox}
    \item \textit{Báo cáo, thảo luận:} 1 học sinh giải thích miệng, lớp nhận xét.
    \item \textit{Kết luận, nhận định:} Giáo viên chuẩn hóa kiến thức.
\end{itemize}

\paragraph{Nội dung 2: Đồ thị và sự biến thiên của 4 hàm số lượng giác}
\textbf{a) Mục tiêu:} Vẽ và mô tả được hình dạng đồ thị của $y = \sin x, y = \cos x, y = \tan x, y = \cot x$; giải thích được các khoảng đồng biến, nghịch biến dựa vào đồ thị; thực hành phòng máy với GeoGebra.

\textbf{b) Nội dung:}
\begin{itemize}[leftmargin=0.6cm]
    \item \textbf{Đồ thị hàm số $y = \sin x$ (Đường hình sin):}
    \begin{itemize}[leftmargin=0.5cm]
        \item Lập bảng giá trị trên một chu kì $[-\pi; \pi]$.
        \item Đồng biến trên mỗi khoảng $\left(-\dfrac{\pi}{2} + k2\pi; \dfrac{\pi}{2} + k2\pi\right)$; nghịch biến trên mỗi khoảng $\left(\dfrac{\pi}{2} + k2\pi; \dfrac{3\pi}{2} + k2\pi\right)$.
        \item Đồ thị đi qua gốc tọa độ $O(0; 0)$ và nhận $O$ làm tâm đối xứng.
    \end{itemize}
    \item \textbf{Đồ thị hàm số $y = \cos x$ (Đường hình sin dịch chuyển):}
    \begin{itemize}[leftmargin=0.5cm]
        \item Lấy được từ đồ thị $y = \sin x$ bằng cách tịnh tiến sang trái một đoạn $\dfrac{\pi}{2}$ (vì $\cos x = \sin\left(x + \dfrac{\pi}{2}\right)$).
        \item Đồng biến trên $(-\pi + k2\pi; k2\pi)$; nghịch biến trên $(k2\pi; \pi + k2\pi)$.
        \item Đồ thị nhận trục tung $Oy$ làm trục đối xứng.
    \end{itemize}
    \item \textbf{Đồ thị hàm số $y = \tan x$:}
    \begin{itemize}[leftmargin=0.5cm]
        \item Luôn \textbf{đồng biến} trên mỗi khoảng xác định $\left(-\dfrac{\pi}{2} + k\pi; \dfrac{\pi}{2} + k\pi\right)$.
        \item Nhận các đường thẳng $x = \dfrac{\pi}{2} + k\pi$ làm tiệm cận đứng.
    \end{itemize}
    \item \textbf{Đồ thị hàm số $y = \cot x$:}
    \begin{itemize}[leftmargin=0.5cm]
        \item Luôn \textbf{nghịch biến} trên mỗi khoảng xác định $(k\pi; \pi + k\pi)$.
        \item Nhận các đường thẳng $x = k\pi$ làm tiệm cận đứng.
    \end{itemize}
\end{itemize}

\textbf{c) Sản phẩm:} Học sinh vẽ phác thảo đúng hình dạng 4 đồ thị vào vở kẻ ô li và hiển thị đồ thị hoàn chỉnh trên GeoGebra.

\textbf{d) Tổ chức thực hiện (Tích hợp Năng lực số 3.1NC1b & Prompt AI 11.C3.1):}
\begin{itemize}[leftmargin=0.6cm]
    \item \textit{Chuyển giao nhiệm vụ:} Học sinh mở GeoGebra nhập lệnh: \texttt{f(x) = sin(x)}, \texttt{g(x) = cos(x)}, \texttt{h(x) = tan(x)}, \texttt{p(x) = 1/tan(x)} để quan sát dạng đồ thị.
    \item \textit{Thực hiện nhiệm vụ (Năng lực AI Mã 11.C3.1):}
    \begin{aibox}{Thực hành Năng lực AI: Phân tích khoảng biến thiên hàm số}
        \textbf{Kỹ thuật Prompt hỏi AI:}\\
        \texttt{"Hãy lập bảng biến thiên và chỉ rõ các khoảng đồng biến, nghịch biến của hàm số $y = \cos x$ trên đoạn $[-\pi; \pi]$. Nêu rõ các điểm cực đại và cực tiểu."}\\
        $\implies$ Học sinh đối chiếu phản hồi của AI:
        \begin{itemize}[leftmargin=0.5cm]
            \item Trên $[-\pi; 0]$: hàm số đồng biến từ $-1$ lên $1$. Điểm cực đại tại $x = 0$ ($y = 1$).
            \item Trên $[0; \pi]$: hàm số nghịch biến từ $1$ xuống $-1$. Các điểm cực tiểu tại $x = \pm \pi$ ($y = -1$).
        \end{itemize}
    \end{aibox}
    \item \textit{Báo cáo, thảo luận:} Đại diện học sinh chỉ trên đồ thị GeoGebra các đỉnh sóng (cực đại) và đáy sóng (cực tiểu).
    \item \textit{Kết luận, nhận định:} Giáo viên chốt các tính chất then chốt: hàm $\tan$ luôn đồng biến trên từng khoảng xác định, hàm $\cot$ luôn nghịch biến.
\end{itemize}

\subsubsection*{3. Hoạt động 3: Luyện tập}
\textbf{a) Mục tiêu:} Rèn luyện kỹ năng đọc đồ thị; tìm tập xác định của hàm lượng giác phức hợp; xác định tính đồng biến, nghịch biến.

\textbf{b) Nội dung:} Học sinh hoàn thành Phiếu học tập số 2 (Worksheet 2):
\begin{itemize}[leftmargin=0.6cm]
    \item \textbf{Bài 1:} Tìm tập xác định của các hàm số sau:
    $$a) \ y = \tan\left(2x - \dfrac{\pi}{3}\right); \quad b) \ y = \dfrac{\cot x}{\cos x - 1}$$
    \item \textbf{Bài 2:} Dựa vào đồ thị hàm số $y = \sin x$, hãy tìm tất cả các giá trị của $x \in [-\pi; 2\pi]$ để $\sin x > 0$.
    \item \textbf{Bài 3:} Xét tính đơn điệu (đồng biến hay nghịch biến) của hàm số $y = \cos x$ trên khoảng $\left(\dfrac{11\pi}{4}; \dfrac{13\pi}{4}\right)$.
\end{itemize}

\textbf{c) Sản phẩm:} Lời giải chi tiết:
\begin{itemize}[leftmargin=0.6cm]
    \item \textbf{Bài 1:}
    \begin{itemize}[leftmargin=0.5cm]
        \item a) Hàm số xác định $\iff 2x - \dfrac{\pi}{3} \neq \dfrac{\pi}{2} + k\pi \iff 2x \neq \dfrac{5\pi}{6} + k\pi \iff x \neq \dfrac{5\pi}{12} + k\dfrac{\pi}{2}$ ($k \in \mathbb{Z}$).\\
        Vậy $D = \mathbb{R} \setminus \left\{\dfrac{5\pi}{12} + k\dfrac{\pi}{2}, k \in \mathbb{Z}\right\}$.
        \item b) Hàm số xác định $\iff \begin{cases} \sin x \neq 0 \ (\text{điều kiện của } \cot) \\ \cos x - 1 \neq 0 \end{cases} \iff \begin{cases} x \neq k\pi \\ \cos x \neq 1 \iff x \neq k2\pi \end{cases} \iff x \neq k\pi$ ($k \in \mathbb{Z}$).\\
        Vậy $D = \mathbb{R} \setminus \{k\pi, k \in \mathbb{Z}\}$.
    \end{itemize}
    \item \textbf{Bài 2:}
    Quan sát đồ thị $y = \sin x$ trên $[-\pi; 2\pi]$, phần đồ thị nằm phía trên trục hoành ($y > 0$) ứng với:
    $$x \in (0; \pi)$$
    (Lưu ý: trên $[-\pi; 0]$ và $[\pi; 2\pi]$ đồ thị nằm phía dưới hoặc cắt trục hoành).
    \item \textbf{Bài 3:}
    Ta có: $\dfrac{11\pi}{4} = 2\pi + \dfrac{3\pi}{4}$ và $\dfrac{13\pi}{4} = 2\pi + \dfrac{5\pi}{4}$.\\
    Do tính tuần hoàn chu kì $2\pi$, sự biến thiên trên $\left(\dfrac{11\pi}{4}; \dfrac{13\pi}{4}\right)$ giống như trên $\left(\dfrac{3\pi}{4}; \dfrac{5\pi}{4}\right)$.\\
    Khoảng này chứa điểm $\pi$.
    \begin{itemize}[leftmargin=0.5cm]
        \item Trên $\left(\dfrac{3\pi}{4}; \pi\right) \subset (0; \pi) \implies$ hàm số \textbf{nghịch biến}.
        \item Trên $\left(\pi; \dfrac{5\pi}{4}\right) \subset (\pi; 2\pi) \implies$ hàm số \textbf{đồng biến}.
    \end{itemize}
    Vậy hàm số không đơn điệu trên toàn khoảng mà nghịch biến trên $\left(\dfrac{11\pi}{4}; 3\pi\right)$ và đồng biến trên $\left(3\pi; \dfrac{13\pi}{4}\right)$.
\end{itemize}

\textbf{d) Tổ chức thực hiện:}
\begin{itemize}[leftmargin=0.6cm]
    \item \textit{Chuyển giao nhiệm vụ:} Học sinh làm bài độc lập trong 8 phút.
    \item \textit{Thực hiện nhiệm vụ (Giàn giáo cho HS TB-Yếu):} Hướng dẫn học sinh yếu cách giải $2x - \dfrac{\pi}{3} \neq \dfrac{\pi}{2} + k\pi$ từng bước: chuyển vế cộng phân số rồi mới chia 2 cho cả hai vế (kể cả số hạng $k\pi$).
    \item \textit{Báo cáo, thảo luận:} 3 học sinh lên bảng trình bày, cả lớp đối chiếu.
    \item \textit{Kết luận, nhận định:} Giáo viên nhận xét, nhấn mạnh lỗi sai thường gặp khi chia hệ số $2$ quên chia cho đuôi chu kì $k\pi$.
\end{itemize}

\subsubsection*{4. Hoạt động 4: Vận dụng}
\textbf{a) Mục tiêu:} Vận dụng hàm số lượng giác giải bài toán thực tế dự báo mực nước thủy triều phục vụ tàu thuyền cập cảng.

\textbf{b) Nội dung:}
Độ sâu $h$ (tính bằng mét) của mực nước tại một cảng biển theo thời gian $t$ (giờ trong ngày, $0 \le t < 24$) được mô hình hóa bởi hàm số lượng giác:
$$h(t) = 3\cos\left(\dfrac{\pi t}{6}\right) + 8$$
\begin{enumerate}[leftmargin=0.5cm]
    \item Độ sâu lớn nhất và nhỏ nhất của mực nước cảng biển là bao nhiêu mét?
    \item Một tàu chở hàng cần độ sâu tối thiểu là $9{,}5\text{ m}$ để có thể cập cảng an toàn. Hỏi vào những thời điểm nào trong ngày con tàu có thể cập cảng?
\end{enumerate}

\textbf{c) Sản phẩm:} Lời giải:
\begin{enumerate}[leftmargin=0.5cm]
    \item Vì $-1 \le \cos\left(\dfrac{\pi t}{6}\right) \le 1 \implies -3 + 8 \le h(t) \le 3 + 8 \iff 5 \le h(t) \le 11$.\\
    Vậy độ sâu lớn nhất là $11\text{ m}$ (khi thủy triều đạt đỉnh) và nhỏ nhất là $5\text{ m}$.
    \item Để tàu cập cảng an toàn, độ sâu phải thỏa mãn:
    $$h(t) \ge 9{,}5 \iff 3\cos\left(\dfrac{\pi t}{6}\right) + 8 \ge 9{,}5 \iff \cos\left(\dfrac{\pi t}{6}\right) \ge \dfrac{1{,}5}{3} = \dfrac{1}{2}$$
    Vì $\cos \ge \dfrac{1}{2} \iff -\dfrac{\pi}{3} + k2\pi \le \dfrac{\pi t}{6} \le \dfrac{\pi}{3} + k2\pi \iff -2 + 12k \le t \le 2 + 12k$ ($k \in \mathbb{Z}$).\\
    Vì $0 \le t < 24$, xét các giá trị nguyên $k$:
    \begin{itemize}[leftmargin=0.5cm]
        \item Với $k = 0 \implies 0 \le t \le 2$ (từ 0 giờ đến 2 giờ sáng).
        \item Với $k = 1 \implies 10 \le t \le 14$ (từ 10 giờ sáng đến 14 giờ chiều).
        \item Với $k = 2 \implies 22 \le t < 24$ (từ 22 giờ đêm đến 24 giờ đêm).
    \end{itemize}
    Vậy tàu có thể cập cảng trong 3 khoảng thời gian: [0h; 2h], [10h; 14h] và [22h; 24h].
\end{enumerate}

\textbf{d) Tổ chức thực hiện:}
\begin{itemize}[leftmargin=0.6cm]
    \item \textit{Chuyển giao nhiệm vụ:} Giáo viên giao bài toán thực tiễn hàng hải cho các nhóm thảo luận trong 4 phút.
    \item \textit{Thực hiện nhiệm vụ:} Nhóm thảo luận, sử dụng đường tròn lượng giác hoặc đồ thị cosin để tìm nghiệm bất phương trình.
    \item \textit{Báo cáo, thảo luận:} Đại diện 1 nhóm lên bảng trình bày giải pháp.
    \item \textit{Kết luận, nhận định:} Giáo viên tổng kết toàn bài: Đồ thị và hàm số lượng giác là cầu nối toán học quan trọng giải thích các quy luật lặp tuần hoàn trong vũ trụ và kỹ thuật. Giao bài tập về nhà trong SGK.
\end{itemize}

% =========================================================================
% PHỤ LỤC
% =========================================================================
\newpage
\section*{IV. PHỤ LỤC HỌC LIỆU BỔ TRỢ}

\subsection*{1. Hệ thống Phiếu học tập (Worksheet) phân tầng bậc thang}

\begin{pedabox}{PHIẾU HỌC TẬP SỐ 1: HÀM SỐ TUẦN HOÀN, SIN VÀ COSIN (Phân tầng 3 mức độ)}
\textbf{Họ và tên học sinh:} \dotfill \textbf{Lớp:} 11A\dots \ \textbf{Ngày:} \dotfill

\textbf{MỨC 1: NHẬN BIẾT (Điền khuyết có giàn giáo)}
\begin{enumerate}[leftmargin=0.5cm]
    \item Điền vào chỗ trống:\\
    a) Đồ thị hàm số chẵn nhận trục \dotfill làm trục đối xứng.\\
    b) Đồ thị hàm số lẻ nhận gốc \dotfill làm tâm đối xứng.\\
    c) Hàm số $y = \sin x$ và $y = \cos x$ đều có tập giá trị là đoạn $[\dots\dots; \dots\dots]$ và tuần hoàn với chu kì $T = \dots\dots$.
    \item Điền từ "chẵn" hoặc "lẻ": Hàm số $y = \cos x$ là hàm số \dotfill; hàm số $y = \sin x$ là hàm số \dotfill.
\end{enumerate}

\textbf{MỨC 2: THÔNG HIỂU (Tìm tập xác định và giá trị lớn nhất, nhỏ nhất)}
\begin{enumerate}[leftmargin=0.5cm]
    \item Tìm tập xác định của hàm số: $y = \dfrac{2}{\cos x}$.\\
    \textit{Giàn giáo gợi ý:} Điều kiện mẫu số khác 0: $\cos x \neq 0 \iff x \neq \dfrac{\pi}{2} + k\pi$.
    \item Tìm giá trị lớn nhất và nhỏ nhất của hàm số: $y = 2\sin x + 5$.
\end{enumerate}

\textbf{MỨC 3: VẬN DỤNG (Khảo sát tính chẵn lẻ)}
\begin{enumerate}[leftmargin=0.5cm]
    \item Xét tính chẵn, lẻ của hàm số: $f(x) = \dfrac{\sin^3 x}{x^2 + 1}$.
\end{enumerate}
\end{pedabox}

\vspace{8pt}

\begin{pedabox}{PHIẾU HỌC TẬP SỐ 2: HÀM SỐ TANG, COTANG VÀ ĐỒ THỊ LƯỢNG GIÁC}
\textbf{MỨC 1: NHẬN BIẾT}
\begin{enumerate}[leftmargin=0.5cm]
    \item Điền vào chỗ trống:\\
    a) Hàm số $y = \tan x$ tuần hoàn với chu kì $T = \dots\dots$ và luôn \dotfill trên mỗi khoảng xác định.\\
    b) Hàm số $y = \cot x$ có tập xác định $D = \mathbb{R} \setminus \{\dots\dots\dots\dots, k \in \mathbb{Z}\}$.
    \item Đúng (Đ) hay Sai (S):\\
    a) Tập giá trị của hàm số $y = \tan x$ là đoạn $[-1; 1]$. (\dots)\\
    b) Đồ thị hàm số $y = \tan x$ nhận gốc tọa độ $O$ làm tâm đối xứng. (\dots)
\end{enumerate}

\textbf{MỨC 2: THÔNG HIỂU}
\begin{enumerate}[leftmargin=0.5cm]
    \item Tìm tập xác định của hàm số: $y = \tan\left(x + \dfrac{\pi}{4}\right)$.\\
    \textit{Giàn giáo:} $x + \dfrac{\pi}{4} \neq \dfrac{\pi}{2} + k\pi \iff x \neq \dots\dots$
    \item Dựa vào đồ thị hàm số $y = \cos x$, hãy tìm các giá trị của $x \in [0; 2\pi]$ sao cho $\cos x = 0$.
\end{enumerate}

\textbf{MỨC 3: VẬN DỤNG}
\begin{enumerate}[leftmargin=0.5cm]
    \item Tìm giá trị lớn nhất và nhỏ nhất của hàm số: $y = \cos^2 x + 2\sin^2 x - 3$.\\
    \textit{Giàn giáo:} Biến đổi $\cos^2 x = 1 - \sin^2 x$ để đưa hàm số về biểu thức bậc nhất theo $\sin^2 x$.
\end{enumerate}
\end{pedabox}

\subsection*{2. Rubric đánh giá năng lực học sinh theo 4 mức độ}

\begin{table}[h!]
\centering
\small
\begin{tabularx}{\textwidth}{|p{2.8cm}|X|X|X|X|}
\hline
\textbf{Tiêu chí} & \textbf{Chưa đạt (Mức 1)} & \textbf{Đạt (Mức 2)} & \textbf{Khá (Mức 3)} & \textbf{Tốt (Mức 4)} \\
\hline
\textbf{1. Khái niệm \& Chu kỳ tuần hoàn} & Nhầm lẫn chu kỳ của $\sin, \cos$ với $\tan, \cot$; chưa nhớ định nghĩa hàm chẵn/lẻ. & Nhớ đúng chu kỳ $2\pi$ và $\pi$; nhận biết được trục đối xứng $Oy$ và tâm đối xứng $O$. & Tìm được chu kỳ của hàm hợp $y = \sin(\omega x)$; xét đúng tính chẵn lẻ của hàm số phức hợp. & Chứng minh chặt chẽ tính tuần hoàn của hàm số bất kỳ; giải thích nguồn gốc chu kỳ qua đường tròn lượng giác. \\
\hline
\textbf{2. Tập xác định \& Tập giá trị} & Quên điều kiện mẫu số của hàm $\tan, \cot$; nhầm tập giá trị của $\tan$ với $[-1; 1]$. & Tìm đúng điều kiện mẫu số cơ bản $x \neq k\pi$; nhớ tập giá trị của $\sin, \cos$ là $[-1; 1]$. & Tìm đúng tập xác định của hàm hợp dạng $\tan(ax + b)$; tìm được $\max, \min$ hàm bậc nhất lượng giác. & Giải quyết trọn vẹn bài toán tìm tập giá trị của hàm lượng giác chứa căn thức hoặc phân thức quy về bậc hai. \\
\hline
\textbf{3. Đọc \& Vẽ đồ thị} & Chưa nhận diện được hình dạng sóng hình sin; vẽ sai các khoảng đồng biến, nghịch biến. & Vẽ phác thảo được đồ thị trên một chu kì; chỉ ra được các điểm đặc biệt cắt trục tọa độ. & Đọc thành thạo đồ thị để giải bất phương trình lượng giác cơ bản; xác định đúng khoảng đơn điệu. & Vẽ đồ thị chính xác trên GeoGebra; phân tích sâu sắc mối liên hệ tịnh tiến giữa đồ thị sin và cosin. \\
\hline
\textbf{4. Năng lực số \& Năng lực AI} & Chưa thao tác được GeoGebra; chưa biết đặt prompt phân tích hàm số. & Mở được đồ thị trên GeoGebra; dùng lệnh vẽ hàm số cơ bản. & Nhập đúng prompt cho AI phân tích bảng biến thiên; dùng bảng tính \texttt{TABLE} khảo sát chu kỳ. & Phân tích sâu sắc ứng dụng AI nhận dạng dữ liệu chu kỳ y tế (ECG); đối chiếu, phát hiện hạn chế của AI. \\
\hline
\end{tabularx}
\end{table}

\end{document}